% Nouveau EGA V (cote 25): archive cote 25, site page 28 (PDF page 29).
% typescript. Draft compared: none.
% AI-transcribed by Claude Opus 5.5 (claude-opus-5-5, adaptive thinking, effort high), Message Batches API, batch ega25b (msgbatch_01A7n5YuXubUPNK1zsepi8RG), 2026-11-13; original AI transcription, subsequently checked against both supplied scan parts; see PAGES_CHECKED.tsv and CHECK_DECISIONS.tsv.
\scanpage{25/28}
\mspage{\hand{29}}
pour tout t \hand{$\in$}U,  sa fibre \var{\inserted{$Y_{\xi}(t)$}}{$Y_{\xi}(t)$}{insertion}{ega25_25_028-v0001} en  t satisfasse la propriété P.  Lorsqu'on ne
dispose pas à priori de \hand{$\xi$}$_{o}$ , \var{\xout{xxx}}{}{typewriter strike-out}{ega25_25_028-v0002} mais que k(s) est infini, on pourra combi-%
ner les deux remarques précédentes pour obtenir une section hyperplane
sur un voisinage ouvert de s, ayant la propriété précédente. \var{\xout{Dixxxxx}}{}{typewriter strike-out}{ega25_25_028-v0003} Enfin,
utilisant \var{\deleted{4.8.}}{}{deletion}{ega25_25_028-v0004}\var{\inserted{le N$^{o}$5}}{le N$^{o}$5}{insertion}{ega25_25_028-v0005} , on a un critère permettant d'affirmer que (X resp. F étant
supposé plat sur S, ce qui permet d'appliquer \var{\deleted{4.7}}{}{deletion}{ega25_25_028-v0006}\var{\inserted{loc. cit.}}{loc. cit.}{insertion}{ega25_25_028-v0007}.)  Y\hand{$_{\xi}$} \var{\inserted[typed]{resp. G\hand{$_{\xi}$}}}{resp. G\hand{$_{\xi}$}}{insertion}{ega25_25_028-v0008} est également
plat sur S. On peut alors, remplaçant X par Y\hand{$_{\xi}$} , itèrer la construction
précédente, ce qui permett\var{\xout{xxxxxx}}{}{typewriter strike-out}{ega25_25_028-v0009} \var{\inserted[typed]{par exemple\var{\xout{\unsure{,|x}}}{}{typewriter strike-out}{ega25_25_028-v0011} sous certaines conditions}}{par exemple\var{\xout{\unsure{,|x}}}{}{typewriter strike-out}{ega25_25_028-v0011} sous certaines conditions}{insertion}{ega25_25_028-v0010} de construire de proche en proche une
``multisection'' \var{\inserted[typed]{S'}}{S'}{insertion}{ega25_25_028-v0012} de X au dessus d'un voisinage ouvert U d'un point donné s,
telle\var{\xout{x}}{}{typewriter strike-out}{ega25_25_028-v0013} que S'\hand{$\longrightarrow$}U soit fini, plat et à fibres satisfaisant la propriété P.
Lorsque k(s) est fini, on peut être contraint de faire un changement de bas%
% NOTE: right edge of the page cut; last letter of "base" not visible
étale surjectif S'\hand{$\longrightarrow$} U (U voisinage ouvert de s) avant de pouvoir appliquer
les constructions précédentes; en effet, sous les conditions du début de
6.1. , si k est fini, il n'existe pas nécessairement de point rationnel sur
k dans  l'ouvert non vide U, mais \var{\xout{l'ensemble des points fermés de U est}}{}{typewriter strike-out}{ega25_25_028-v0014}
il existe certainement un point fermé de U, donc un point à valeurs dans
une extension f inie k' (nécessairement séparable) de k; lorsque k = k(s),
on peut donc, quitte à faire une extension \var{\inserted[typed]{finie}}{finie}{insertion}{ega25_25_028-v0015} étale convenable \var{\inserted[typed]{S'}}{S'}{insertion}{ega25_25_028-v0016} sur un voisinag%
% NOTE: right edge cut; end of "voisinage" not visible
ouvert U de  s, correspondant à l'extension résiduelle k'/k\var{\xout{xxxxxxxxx}}{}{typewriter strike-out}{ega25_25_028-v0017}
i.e. telle que  S$'_{s}$ \hand{$\xrightarrow{\sim}$}\var{\inserted{k(s)-isomorphisme}}{k(s)-isomorphisme }{insertion}{ega25_25_028-v0018}Spec(k\hand{'}) , se ramener à la situation favorable \var{\xout{par}}{}{typewriter strike-out}{ega25_25_028-v0019}
en l'unique point s'\hand{$\in$}S' au dessus de s, par le changement de base S'\hand{$\to$} S.

Donnons quelques exemples d'énoncés obtenus par ce\var{\over{s}{tte}}{tte}{writing over}{ega25_25_028-v0020} méthode\var{\deleted{s}}{}{deletion}{ega25_25_028-v0021}. \var{\deleted{Il y a}}{}{deletion}{ega25_25_028-v0022}
\var{\deleted{lieu cependant, de faire une remarque préliminaire, pour l'élimination dan}}{}{deletion}{ega25_25_028-v0023}%
% NOTE: right edge cut ("dans")
\var{\deleted{\uline{\var{\xout{Propositionxxxxxxxxxx}}{}{typewriter strike-out}{ega25_25_028-v0025}}}}{}{deletion}{ega25_25_028-v0024}
\var{\deleted{ces constructions de l'hypothèse faite dans 4.6. \var{\deleted{et}}{}{deletion}{ega25_25_028-v0027} 4.7. \var{\inserted{et 4.8.}}{et 4.8.}{insertion}{ega25_25_028-v0028} que  aucune compo}}{}{deletion}{ega25_25_028-v0026}%
\var{\deleted{sante irréductible géométrique d\var{\xout{x}}{}{typewriter strike-out}{ega25_25_028-v0030}'un supp F$_{s}$ n'est contenue\hand{,}dans une sectio}}{}{deletion}{ega25_25_028-v0029}%
\var{\deleted{hyperplane (definie sur $\hand{\overline{\textcolor{black}{\text{k(s)}}}}$ ): on peut (et doit) en effet énoncer 4.6.}}{}{deletion}{ega25_25_028-v0031}
\var{\deleted{et 4.7. sous forme légèrement plus générale, en se donnant une partie ouver}}{}{deletion}{ega25_25_028-v0032}%
\var{\deleted{te W de $\overset{\hand{\vee}}{P}$ , \var{\xout{ex}}{}{typewriter strike-out}{ega25_25_028-v0034}  en remplaçant dans la condition envisagée  ``hyperplan'' par}}{}{deletion}{ega25_25_028-v0033}
\var{\margin{28}}{\placedmargin{28}}{marginal note placed}{ega25_25_028-v0035}
